\newpage
\titledquestion{Inversion Preserves the Verdict}

\textbf{Inversion Preserves the Verdict}

Recall the type of polymorphic trees:
\begin{codeblock}
  datatype 'a tree = Empty | Node of 'a tree * 'a * 'a tree
\end{codeblock}

and the \code{invert} function, which mirrors a tree by swapping every left and
right subtree:
\begin{codeblock}
  fun invert Empty = Empty
    | invert (Node (L, x, R)) = Node (invert R, x, invert L)
\end{codeblock}

We now define \code{treeAll}, which checks whether \textbf{every} element of a
tree satisfies a predicate:
\begin{codeblock}
  fun treeAll p Empty = true
    | treeAll p (Node (L, x, R)) =
        p x andalso treeAll p L andalso treeAll p R
\end{codeblock}

Intuitively, mirroring a tree should not change \textit{which} elements it
contains --- only their arrangement. So whether all of them satisfy \code{p}
should be unaffected by inversion.

\begin{parts}
  \part[18]\hfill

    Prove the following theorem.

    \textbf{Theorem.} For all values \code{p : 'a -> bool} and all values
    \code{T : 'a tree}:
    $$\code{treeAll p (invert T)} \eeq \code{treeAll p T}$$

    You may use the following lemmas without proof:

    \textbf{Lemma 1:} \code{invert} is total.

    \textbf{Lemma 2:} For any total \code{p}, \code{treeAll p} is total.

    \textbf{Lemma 3 (commuting a conjunction):} For all expressions
    \code{a}, \code{b}, \code{c} of type \code{bool} that each reduce to a
    value, \\
    \code{a andalso b andalso c} $\eeq$ \code{a andalso c andalso b}.

    \begin{constraint}
      You may assume \code{p} is total.
    \end{constraint}

    \begin{solutionorbox}[40em]

      By structural induction on \code{T : 'a tree}.

      \vspace{6pt}
      \textbf{Base case:} \code{T = Empty}.

      WTS: \code{treeAll p (invert Empty)} $\eeq$ \code{treeAll p Empty}.

      \begin{align*}
        & \code{treeAll p (invert Empty)} \\
        &\eeq \code{treeAll p Empty} \tag{clause 1 of \code{invert}} \\
        &\eeq \code{true} \tag{clause 1 of \code{treeAll}}
      \end{align*}

      and directly \code{treeAll p Empty} $\eeq$ \code{true} by clause 1 of
      \code{treeAll}, so both sides are equal.

      \vspace{6pt}
      \textbf{Inductive case:} \code{T = Node (L, x, R)}.

      IH1: \code{treeAll p (invert L)} $\eeq$ \code{treeAll p L}

      IH2: \code{treeAll p (invert R)} $\eeq$ \code{treeAll p R}

      WTS: \code{treeAll p (invert (Node (L, x, R)))} $\eeq$
           \code{treeAll p (Node (L, x, R))}

      \vspace{4pt}
      Starting from the left-hand side:
      \begin{align*}
        & \code{treeAll p (invert (Node (L, x, R)))} \\
        &\eeq \code{treeAll p (Node (invert R, x, invert L))}
          \tag{clause 2 of \code{invert}} \\
        &\eeq \code{p x andalso treeAll p (invert R)} \\
        & \qquad \code{andalso treeAll p (invert L)}
          \tag{clause 2 of \code{treeAll}, Lemma 1} \\
        &\eeq \code{p x andalso treeAll p R andalso treeAll p L}
          \tag{IH1 and IH2} \\
        &\eeq \code{p x andalso treeAll p L andalso treeAll p R}
          \tag{Lemma 3, Lemma 2} \\
        &\eeq \code{treeAll p (Node (L, x, R))}
          \tag{clause 2 of \code{treeAll}}
      \end{align*}

      which is the right-hand side, completing the inductive case.

      \vspace{4pt}
      By the principle of structural induction, the theorem holds for all
      \code{T : 'a tree}. $\qed$

    \end{solutionorbox}
\end{parts}
